\documentclass[a4paper,10pt]{article}
\usepackage[utf8]{inputenc}
\usepackage{graphicx}
\usepackage{listings}
\usepackage{caption}
\usepackage{subcaption}
\usepackage{enumerate}

\usepackage{listings}
\usepackage[usenames,dvipsnames]{color}
\usepackage{courier}


% Title Page
\title{}
\author{}


\begin{document}

\begin{titlepage}
\begin{center}


\textsc{\LARGE Royal Institute of Technology}\\[1.5cm]


%\includegraphics[width=0.3\textwidth]{kth_mathematics_rgb.jpg}\\[1cm]
\begin{figure}[h]
        \centering
        \begin{subfigure}[b]{0.2\textwidth}
                \centering
                \includegraphics[width=\textwidth]{kth_mathematics_rgb.jpg}
                \caption{Caroline}
                \label{fig:Caroline}
        \end{subfigure}%
        ~ %add desired spacing between images, e. g. ~, \quad, \qquad etc. 
          %(or a blank line to force the subfigure onto a new line)
        \begin{subfigure}[b]{0.2\textwidth}
                \centering
                \includegraphics[width=\textwidth]{kth_mathematics_rgb.jpg}
                \caption{Paul}
                \label{fig:Paul}
        \end{subfigure}
        \\ %add desired spacing between images, e. g. ~, \quad, \qquad etc. 
          %(or a blank line to force the subfigure onto a new line)
        \begin{subfigure}[b]{0.2\textwidth}
                \centering
                \includegraphics[width=\textwidth]{kth_mathematics_rgb.jpg}
                \caption{Sindri}
                \label{fig:Sindri}
        \end{subfigure}
          ~ %add desired spacing between images, e. g. ~, \quad, \qquad etc. 
          %(or a blank line to force the subfigure onto a new line)
        \begin{subfigure}[b]{0.2\textwidth}
                \centering
                \includegraphics[width=\textwidth]{kth_mathematics_rgb.jpg}
                \caption{Tobias}
                \label{fig:Tobias}
        \end{subfigure}
        %\caption{Pictures of animals}\label{fig:animals}
\end{figure}

\textsc{\Large Artificial intelligence, dd2380 }  %\\ Instructor: Tobias Ryden , \emph{tryd@math.kth.se}}\\[0.5cm]

\hrulefill \\[0.4cm]
{ \huge \bfseries Final project: sokoban}\\[0.4cm]
\hrulefill \\[1.5cm]



    Caroline Laurène Kéramsi, keramsi@kth.se, XXXXXX-XXXX \\
    Paul Lagrée,lagree@kth.se, XXXXXX-XXXX \\
    Sindri Magnússon, sindrim@kth.se, 871209-7156 \\
    Tobias Johannes, Uebbing@kth.se, XXXXXX-XXXX \\


\vfill
% Bottom of the page
{\large \today}

\end{center}
\end{titlepage}


%%%%%%%%%%%%%%%%%%%%%%%%%%%%%%%%%%%%%%%%%%%%%%%%%%%%%%%%%%%%%%%%%%%%%%%%%%%%%%%%
%%%%%%%%%%%%%%%%%%%%%%%%%%%%%%%%%%%%%%%%%%%%%%%%%%%%%%%%%%%%%%%%%%%%%%%%%%%%%%%%
%%%%%%%%%%%%%%%%%%%%%%%%%%%%%%%%%%%%%%%%%%%%%%%%%%%%%%%%%%%%%%%%%%%%%%%%%%%%%%%%
%%%%%%%%%%
%%%%%%%%%% Report statrs
%%%%%%%%%%
\cleardoublepage
\tableofcontents
\newpage

\abstract{ asdf}



\section{Introduction}

    Sokoban is an popular puzzel dating back to the eighties.  
    The original setting of the puzzel is a warehouse and the problem is to push boxes
    around to predefined storage locations.  
    But the underlying problem is of course much more abstract.  The rules are simple~\cite{wiki_soko}: 
    %cite http://en.wikipedia.org/wiki/Sokoban
    \begin{enumerate}[I]
        \item  Only one box can be pushed at a time.
        \item A box cannot be pulled.
        \item The player cannot walk through boxes or walls.
        \item The puzzle is solved when all boxes are located at storage locations.
    \end{enumerate}
    Though the rules are simple the problem is quite difficult and has been proven to be $NP$-hard.
    This is not only due to the brancing factor but also the 
    enormous deepth of the search tree~\cite{wiki_soko}.
    

    
   

\section{Design}

    \subsection{States}
    \subsection{Search}
        sdf~\cite{solving_soko}
    \subsection{Pruning}


\section{Results}

\section{Conclusions}
   
      
%%%%%%%%%%%%%%%%%%%%%%%%%%%%%%%%%%%%%%%%%%
%%%%%%%%%%%%%%%%%%%%%%%%%%%%%%%%%%%%%%%%%%
%%%%%%
%%%%%%  Appendix  
%%%%%%
%\appendix
 
%%%%%%%%%%%%%%%%%%%%%%%%%%%%%%%%%%%%%%%%%%
%%%%%%%%%%%%%%%%%%%%%%%%%%%%%%%%%%%%%%%%%%
%%%%%%
%%%%%%    
%%%%%% 
\bibliographystyle{plain} 
\bibliography{refs.bib}	
\end{document}          
